\documentclass[11pt,a4paper]{article}

% --- Packages ---
\usepackage[utf8]{inputenc}
\usepackage[T1]{fontenc}
\usepackage{amsmath,amssymb,amsthm}
\usepackage{physics}
\usepackage{geometry}
\usepackage{enumitem}
\usepackage{hyperref}
\usepackage{titlesec}
\usepackage{fancyhdr}
\usepackage{mathtools}

\geometry{margin=1in}
\pagestyle{fancy}
\fancyhf{}
\rhead{Quantum Spin Lecture Notes}
\lhead{Dirac's Belt Trick, Topology, and Spin-\textonehalf{} Particles}
\cfoot{\thepage}

% --- Theorem Environments ---
\theoremstyle{definition}
\newtheorem{definition}{Definition}[section]
\newtheorem{example}{Example}[section]
\theoremstyle{plain}
\newtheorem{theorem}{Theorem}[section]
\newtheorem{proposition}{Proposition}[section]
\newtheorem{corollary}{Corollary}[section]
\theoremstyle{remark}
\newtheorem{remark}{Remark}[section]

\title{\textbf{Dirac's Belt Trick, Topology, and Spin-$\frac{1}{2}$ Particles}}
\author{Lecture Notes on Quantum Spin}
\date{}

\begin{document}
\maketitle
\tableofcontents
\newpage

% ============================================================
\section{Introduction: The Dirac Belt Trick}
% ============================================================

The \textbf{Dirac belt trick} is a striking physical demonstration invented by Paul Dirac to illustrate the bizarre rotational properties of spin-$\frac{1}{2}$ particles such as electrons, protons, and neutrons.

\subsection{The Setup}

Take a belt with one end fixed (e.g., held in place by a heavy book) and the other end free. The rules are:
\begin{enumerate}[label=(\roman*)]
    \item You are \emph{not} allowed to rotate either end of the belt.
    \item You \emph{are} allowed to translate (move) both ends through space without turning them.
\end{enumerate}

\subsection{The Observation}

\begin{itemize}
    \item \textbf{$360^\circ$ twist:} Twist the free end by $360^\circ$ ($2\pi$ radians). No matter what manipulations you perform (without rotating the ends), \emph{you cannot undo the twist}.
    \item \textbf{$720^\circ$ twist:} Twist the free end by $720^\circ$ ($4\pi$ radians). Remarkably, the twist \emph{can} be undone by looping the belt around, all without rotating either end.
\end{itemize}

This mirrors the behavior of spin-$\frac{1}{2}$ particles: a $2\pi$ rotation does \emph{not} return the particle to its original quantum state, but a $4\pi$ rotation does.

% ============================================================
\section{The Space of Rotations: $\mathrm{SO}(3)$}
% ============================================================

\subsection{Specifying a Rotation}

Any rotation in three-dimensional space is specified by two pieces of information:
\begin{enumerate}
    \item An \textbf{axis of rotation} $\hat{n} \in \mathbb{R}^3$, $|\hat{n}| = 1$.
    \item An \textbf{angle of rotation} $\theta \in [0, \pi]$ (measured in radians, counterclockwise via the right-hand rule).
\end{enumerate}

These can be packaged into a single vector $\vec{v} = \theta\,\hat{n}$, where the direction encodes the axis and the length encodes the angle.

\subsection{Visualizing the Rotation Space}

The space of all rotations can be visualized as a \textbf{solid ball of radius $\pi$} in $\mathbb{R}^3$:
\begin{itemize}
    \item Each point inside the ball corresponds to a unique rotation (given by the vector from the center to that point).
    \item The center of the ball corresponds to the identity rotation (no rotation).
    \item \textbf{Antipodal identification:} A rotation by $\pi$ radians clockwise around an axis is the same as a rotation by $\pi$ radians counterclockwise around the same axis. Therefore, \emph{diametrically opposite points on the boundary of the ball must be identified as the same rotation}.
\end{itemize}

\begin{remark}
    Topologically, $\mathrm{SO}(3)$ is \emph{not} simply a ball; it is a ball with antipodal boundary points identified. This identification is the key to understanding the belt trick.
\end{remark}

\subsection{Rotation Matrices and $\mathrm{SO}(3)$}

Rotations act on position vectors $\vec{r} \in \mathbb{R}^3$ via $3 \times 3$ orthogonal matrices. For example, a rotation about the $x$-axis by angle $\theta$:
\[
    R_x(\theta) = \begin{pmatrix} 1 & 0 & 0 \\ 0 & \cos\theta & -\sin\theta \\ 0 & \sin\theta & \cos\theta \end{pmatrix}.
\]

\begin{definition}[Orthogonal Matrix]
    A matrix $R$ is \textbf{orthogonal} if $R^\top R = I$, i.e., it preserves the inner product:
    \[
        \langle R\vec{r}_1,\, R\vec{r}_2 \rangle = \langle \vec{r}_1,\, \vec{r}_2 \rangle \quad \forall\; \vec{r}_1, \vec{r}_2 \in \mathbb{R}^3.
    \]
\end{definition}

\begin{definition}[$\mathrm{SO}(3)$]
    The \textbf{special orthogonal group} $\mathrm{SO}(3)$ is the group of all $3 \times 3$ real orthogonal matrices with determinant $+1$:
    \[
        \mathrm{SO}(3) = \{ R \in \mathbb{R}^{3\times 3} : R^\top R = I,\; \det R = 1 \}.
    \]
    The ``S'' stands for ``special'' ($\det = 1$), ``O'' for ``orthogonal,'' and ``3'' for three dimensions.
\end{definition}

% ============================================================
\section{Paths, Belts, and Contractibility}
% ============================================================

\subsection{Paths in the Space of Rotations}

A \textbf{path} in $\mathrm{SO}(3)$ is a continuous curve $\gamma: [0,1] \to \mathrm{SO}(3)$. A belt configuration corresponds to such a path: as you move along the belt from one end to the other, the local orientation of coordinate axes traces out a path in rotation space.

\begin{itemize}
    \item A \textbf{perfectly straight belt} (no twist) corresponds to the \textbf{trivial path} --- the path that stays at the identity (center of the ball).
    \item The constraint that neither end of the belt rotates means the path must be a \textbf{loop}: it starts and ends at the identity.
    \item ``Undoing'' a twist means \textbf{contracting} the loop to a point (the trivial path), without moving the endpoints.
\end{itemize}

\subsection{Contractibility Results}

\begin{proposition}
    The loop in $\mathrm{SO}(3)$ corresponding to a $2\pi$ twist is \textbf{not contractible} to the trivial path.
\end{proposition}

\textit{Intuition:} The $2\pi$-twist loop passes through the boundary of the ball, connecting two antipodal points. Since these antipodal points are identified, the path forms a closed loop that ``wraps around'' the space in a way that cannot be undone.

\begin{proposition}
    The loop in $\mathrm{SO}(3)$ corresponding to a $4\pi$ twist \textbf{is contractible} to the trivial path.
\end{proposition}

\textit{Intuition:} The $4\pi$-twist loop passes through the boundary twice. The two pairs of antipodal crossings can be brought together, allowing the path to be pulled through the boundary and contracted.

% ============================================================
\section{Groups}
% ============================================================

\begin{definition}[Group]
    A \textbf{group} is a set $G$ equipped with a binary operation $+$ satisfying:
    \begin{enumerate}[label=(\arabic*)]
        \item \textbf{Closure:} $g_1, g_2 \in G \implies g_1 + g_2 \in G$.
        \item \textbf{Identity:} $\exists\, \mathrm{id} \in G$ such that $g + \mathrm{id} = g$ for all $g \in G$.
        \item \textbf{Associativity:} $(g_1 + g_2) + g_3 = g_1 + (g_2 + g_3)$.
        \item \textbf{Inverses:} For each $g \in G$, $\exists\, (-g) \in G$ such that $g + (-g) = \mathrm{id}$.
    \end{enumerate}
\end{definition}

\begin{example}[Rotations Form a Group]
    The set $\mathrm{SO}(3)$ forms a group under composition of rotations:
    \begin{itemize}
        \item The ``sum'' of two rotations is performing one after the other.
        \item The identity element is the trivial rotation (do nothing).
        \item The inverse of a rotation by $\theta$ about $\hat{n}$ is a rotation by $\theta$ about $-\hat{n}$.
    \end{itemize}
\end{example}

% ============================================================
\section{The Fundamental Group of $\mathrm{SO}(3)$}
% ============================================================

\subsection{Fundamental Group}

\begin{definition}[Fundamental Group $\pi_1(X)$]
    Let $X$ be a topological space with a chosen basepoint $x_0$. The \textbf{fundamental group} $\pi_1(X, x_0)$ is the set of all loops starting and ending at $x_0$, up to continuous deformation (homotopy), with the group operation being concatenation of loops.
\end{definition}

\subsection{$\pi_1(\mathrm{SO}(3)) = \mathbb{Z}/2\mathbb{Z}$}

In $\mathrm{SO}(3)$, with the basepoint at the identity:
\begin{itemize}
    \item The \textbf{trivial loop} $[\mathrm{id}]$: the path that stays at the identity.
    \item The \textbf{non-trivial loop} $[p]$: the path corresponding to the $2\pi$ twist, which is \emph{not} contractible.
\end{itemize}

Adding $[p]$ to itself gives the $4\pi$-twist loop, which \emph{is} contractible:
\[
    [p] + [p] = [\mathrm{id}].
\]

The full group table is:
\[
\begin{array}{c|cc}
    + & [\mathrm{id}] & [p] \\ \hline
    [\mathrm{id}] & [\mathrm{id}] & [p] \\
    [p] & [p] & [\mathrm{id}]
\end{array}
\]

\begin{theorem}
    \[
        \pi_1\bigl(\mathrm{SO}(3)\bigr) \cong \mathbb{Z}/2\mathbb{Z} \cong \mathbb{Z}_2.
    \]
\end{theorem}

The Dirac belt trick is a physical demonstration of this theorem.

% ============================================================
\section{Quantum Spin States and $\mathrm{SU}(2)$}
% ============================================================

\subsection{Spin-$\frac{1}{2}$ States}

Quantum spin states are vectors in $\mathbb{C}^2$:
\[
    \ket{\psi} = \begin{pmatrix} \alpha \\ \beta \end{pmatrix}, \quad \alpha, \beta \in \mathbb{C}.
\]

The inner product on $\mathbb{C}^2$ involves the conjugate transpose:
\[
    \braket{\psi_1}{\psi_2} = \bar{\alpha}_1 \alpha_2 + \bar{\beta}_1 \beta_2.
\]

Physical states satisfy the normalization condition $\braket{\psi}{\psi} = |\alpha|^2 + |\beta|^2 = 1$.

\subsection{Measurement Probabilities}

If an electron has spin state $\ket{\psi} = \begin{pmatrix} \alpha \\ \beta \end{pmatrix}$, then upon measuring spin along the $z$-axis:
\begin{align}
    P(\text{spin up}) &= |\alpha|^2 = \bigl|\braket{\uparrow_z}{\psi}\bigr|^2, \\
    P(\text{spin down}) &= |\beta|^2 = \bigl|\braket{\downarrow_z}{\psi}\bigr|^2.
\end{align}

The spin-up and spin-down states along each axis are:
\[
    \ket{\uparrow_x} = \frac{1}{\sqrt{2}}\begin{pmatrix}1\\1\end{pmatrix}, \quad
    \ket{\downarrow_x} = \frac{1}{\sqrt{2}}\begin{pmatrix}1\\-1\end{pmatrix},
\]
\[
    \ket{\uparrow_y} = \frac{1}{\sqrt{2}}\begin{pmatrix}1\\i\end{pmatrix}, \quad
    \ket{\downarrow_y} = \frac{1}{\sqrt{2}}\begin{pmatrix}1\\-i\end{pmatrix},
\]
\[
    \ket{\uparrow_z} = \begin{pmatrix}1\\0\end{pmatrix}, \quad
    \ket{\downarrow_z} = \begin{pmatrix}0\\1\end{pmatrix}.
\]

\subsection{Rotations of Spin States: $\mathrm{SU}(2)$}

Rotations act on spin states via $2 \times 2$ unitary matrices $U$ with $\det U = 1$.

\begin{definition}[$\mathrm{SU}(2)$]
    The \textbf{special unitary group} $\mathrm{SU}(2)$ is
    \[
        \mathrm{SU}(2) = \{ U \in \mathbb{C}^{2\times 2} : U^\dagger U = I,\; \det U = 1 \}.
    \]
\end{definition}

A general element of $\mathrm{SU}(2)$ has the form
\[
    U = \begin{pmatrix} a & b \\ -\bar{b} & \bar{a} \end{pmatrix}, \quad |a|^2 + |b|^2 = 1.
\]

Writing $a = X + iY$ and $b = Z + iW$, the constraint becomes $X^2 + Y^2 + Z^2 + W^2 = 1$, which is the equation for the \textbf{3-sphere} $S^3 \subset \mathbb{R}^4$:
\[
    \mathrm{SU}(2) \cong S^3 \quad \text{(as topological spaces)}.
\]

% ============================================================
\section{$\mathrm{SU}(2)$ as a Double Cover of $\mathrm{SO}(3)$}
% ============================================================

\subsection{Covering Spaces}

\begin{definition}[Covering Space]
    A \textbf{covering space} of a topological space $X$ is a space $\tilde{X}$ together with a continuous surjective map $p: \tilde{X} \to X$ such that every point of $X$ has a neighborhood evenly covered by $p$. If each fiber $p^{-1}(x)$ has exactly two elements, we call it a \textbf{double cover}.
\end{definition}

\subsection{The Double Cover $\mathrm{SU}(2) \to \mathrm{SO}(3)$}

The key topological relationship is:
\begin{theorem}
    $\mathrm{SU}(2)$ is a double cover of $\mathrm{SO}(3)$: there is a $2$-to-$1$ surjective group homomorphism
    \[
        \phi: \mathrm{SU}(2) \longrightarrow \mathrm{SO}(3).
    \]
    The kernel is $\ker\phi = \{I, -I\}$, so $U$ and $-U$ map to the same rotation.
\end{theorem}

The covering map is constructed using the Pauli matrices. Given $\vec{r} = (x,y,z)$, form the traceless Hermitian matrix
\[
    \vec{r} \cdot \vec{\sigma} = x\sigma_x + y\sigma_y + z\sigma_z = \begin{pmatrix} z & x - iy \\ x + iy & -z \end{pmatrix}.
\]
Under conjugation by $U \in \mathrm{SU}(2)$:
\[
    U(\vec{r}\cdot\vec{\sigma})U^\dagger = \vec{r}' \cdot \vec{\sigma},
\]
where $\vec{r}' = R\vec{r}$ for some $R \in \mathrm{SO}(3)$. This defines the map $\phi(U) = R$.

\subsection{Lifting Paths and Non-Contractibility}

When we lift the $2\pi$-twist loop from $\mathrm{SO}(3)$ to $\mathrm{SU}(2)$:
\begin{itemize}
    \item The lifted path starts at $I \in \mathrm{SU}(2)$ but ends at $-I \in \mathrm{SU}(2)$.
    \item Since the endpoints differ, the lifted path is not a loop, proving by contradiction that the original loop in $\mathrm{SO}(3)$ is non-contractible.
\end{itemize}

The $4\pi$-twist loop lifts to a loop in $\mathrm{SU}(2)$ (starts and ends at $I$), which \emph{is} contractible since $\pi_1(S^3) = \{e\}$ (the trivial group --- every loop on $S^3$ is contractible).

\subsection{Fundamental Group Relationship}

\[
    \pi_1\bigl(\mathrm{SU}(2)\bigr) = \{e\}, \quad \pi_1\bigl(\mathrm{SO}(3)\bigr) = \mathbb{Z}_2.
\]
The fundamental group of $\mathrm{SO}(3)$ has exactly \emph{twice} as many elements as that of $\mathrm{SU}(2)$, reflecting the \emph{double} cover.

% ============================================================
\section{Why Electrons Negate Under $2\pi$ Rotation}
% ============================================================

The physical consequence of all the above mathematics:

\begin{enumerate}
    \item The group that acts on $\mathbb{C}^2$ spin vectors is $\mathrm{SU}(2)$, not $\mathrm{SO}(3)$.
    \item To rotate an electron, we must \textbf{lift} a path from $\mathrm{SO}(3)$ to $\mathrm{SU}(2)$.
    \item A $2\pi$ rotation, when lifted from $I \in \mathrm{SU}(2)$, ends at $-I \in \mathrm{SU}(2)$.
    \item A $4\pi$ rotation returns to $I \in \mathrm{SU}(2)$.
\end{enumerate}

\begin{theorem}[Spin-$\frac{1}{2}$ Rotation Property]
    For a spin-$\frac{1}{2}$ particle:
    \begin{align}
        \text{Rotation by } 2\pi &:\quad \ket{\psi} \mapsto -\ket{\psi}, \\
        \text{Rotation by } 4\pi &:\quad \ket{\psi} \mapsto +\ket{\psi}.
    \end{align}
\end{theorem}

The explicit $\mathrm{SU}(2)$ matrix for a rotation by angle $\theta$ about axis $\hat{n}$ is:
\[
    U(\hat{n}, \theta) = \cos\frac{\theta}{2}\, I - i\sin\frac{\theta}{2}\, (\hat{n}\cdot\vec{\sigma}).
\]
Substituting $\theta = 2\pi$:
\[
    U(\hat{n}, 2\pi) = \cos\pi\, I - i\sin\pi\, (\hat{n}\cdot\vec{\sigma}) = -I.
\]
This confirms that a $2\pi$ rotation negates the spin state.

% ============================================================
\section{Physical Realization and Detection}
% ============================================================

\subsection{How to Rotate an Electron}

\begin{enumerate}
    \item \textbf{Change of reference frame:} Walking $2\pi$ radians around an electron while continuously describing its state in your own frame results in a negated state.
    \item \textbf{Magnetic field (Larmor precession):} An electron's spin angular momentum vector precesses around an external magnetic field. By tuning the field strength and duration, any desired rotation can be achieved.
\end{enumerate}

\subsection{Detecting the Minus Sign}

The overall phase $-1$ from a $2\pi$ rotation is \emph{observable} via interference:
\begin{enumerate}
    \item Place an electron (or neutron) in a quantum superposition at two spatial locations.
    \item Subject one copy to a magnetic field inducing a $2\pi$ rotation (introducing the minus sign).
    \item Recombine the two copies.
\end{enumerate}
Interference effects from the relative minus sign are measurable. This has been confirmed experimentally via \textbf{neutron interferometry}.

% ============================================================
\section{Additional Topological Remarks}
% ============================================================

\subsection{Real Projective Space}

Identifying antipodal points on the 2-sphere $S^2$ yields 2-dimensional real projective space:
\[
    \mathbb{R}P^2 = S^2 / (x \sim -x).
\]
Extending to the 3-sphere:
\[
    \mathbb{R}P^3 = S^3 / (x \sim -x).
\]
Topologically, $\mathrm{SO}(3) \cong \mathbb{R}P^3$, which is consistent with $\mathrm{SU}(2) \cong S^3$ being a double cover.

\subsection{Nature and Spin}

The fundamental group $\pi_1(\mathrm{SO}(3)) = \mathbb{Z}_2$ constrains what kinds of particles nature can produce. A spin-$\frac{1}{3}$ particle is impossible because $\mathbb{Z}_2 \neq \mathbb{Z}_3$. Nature's creativity operates within the rigid bounds of mathematical structure.

\end{document}
